% GENERATED, edit the French unit or the translation table
\mxnewpage\section*{EV4-T2 -- EAM preparation}\label{EV4-T2}
\addcontentsline{toc}{section}{EV4-T2 -- EAM preparation}
\textbf{Practice paper}. Whole programme.\par Duration: 2 h. No calculator. Marked out of 20. Practice sessions in May and June; EAM date for the zone to be confirmed.
\subsection*{Part 1 -- Fluency (6 points)}
Twelve questions, 0.5 point each. Only one correct answer; no justification required. An incorrect or missing answer scores 0, with no loss of points.
\begin{enumerate}
\item Calculate $6+2\times3$.\par A. $24$\quad B. $15$\quad C. $11$\quad D. $12$
\item Calculate $\frac12+\frac13$.\par A. $\frac{5}{6}$\quad B. $\frac{2}{5}$\quad C. $\frac{5}{3}$\quad D. $\frac{1}{6}$
\item Calculate $\dfrac{2^6}{2^{4}}$.\par A. $2$\quad B. $4$\quad C. $1$\quad D. $11$
\item A value of 50 increases by 20 \%. Final value?\par A. $60$\quad B. $10$\quad C. $70$\quad D. $40$
\item A value of 80 is obtained after a 20 \% decrease. Initial value?\par A. $64$\quad B. $100$\quad C. $96$\quad D. $60$
\item A 25 \% increase then a 20 \% decrease: overall coefficient?\par A. $\frac{4}{5}$\quad B. $\frac{21}{20}$\quad C. $1$\quad D. $\frac{9}{10}$
\item The solution of $2x+6=0$ is\par A. $-3$\quad B. $12$\quad C. $3$\quad D. $-12$
\item The slope between $(0;6)$ and $(2;12)$ is\par A. $2$\quad B. $-3$\quad C. $3$\quad D. $6$
\item The mean of $6, 8, 10$ is\par A. $6$\quad B. $24$\quad C. $8$\quad D. $9$
\item A sector of $90^\circ$ represents what fraction of the total?\par A. $\frac{1}{3}$\quad B. $\frac{1}{2}$\quad C. $\frac{3}{4}$\quad D. $\frac{1}{4}$
\item A probability is $\dfrac{2}{5}$. Its complement is\par A. $\frac{5}{2}$\quad B. $\frac{2}{5}$\quad C. $\frac{1}{5}$\quad D. $\frac{3}{5}$
\item For $f(x)=x^2+6$, $f(-2)$ is\par A. $10$\quad B. $8$\quad C. $24$\quad D. $2$
\end{enumerate}
\mxnewpage\subsection*{Part 2 -- Independent exercises (14 points)}
Every step must be justified. Exact results are expected unless stated otherwise.
\mxneed{12}
\subsection*{Exercise 1 -- A reversed tree (5 points)}
Two events $A$ and $B$ satisfy $P(A)=\dfrac35$, $P_A(B)=\dfrac13$ and $P_{\overline A}(B)=\dfrac34$. We wish to build a tree that starts by distinguishing $B$ and $\overline B$, instead of $A$ and $\overline A$.
\begin{enumerate}
\item Calculate $P(A\cap B)$ and $P(B)$. \hfill(1 pt)
\item Calculate $P_B(A)$ and $P_{\overline B}(A)$. \hfill(2 pts)
\item Construct and complete the weighted tree whose first level distinguishes $B$ and $\overline B$, then the second $A$ and $\overline A$. \hfill(2 pts)
\end{enumerate}
\mxneed{12}
\subsection*{Exercise 2 -- Two optimisation methods (5 points)}
The function $f$ is defined on the interval $[1;6]$ by $f(x)=x(8-x)$. We look for its extrema on this precise domain, then for the values of $x$ for which $f(x)$ reaches at least 15.
\begin{enumerate}
\item Give a vertex form and a maximum. \hfill(1 pt)
\item Recover the maximum by differentiation. \hfill(2 pts)
\item Determine the minimum on the domain and solve $f(x)\geq15$. \hfill(2 pts)
\end{enumerate}
\mxneed{12}
\subsection*{Exercise 3 -- Arguments and symmetry (4 points)}
In an orthonormal frame with origin $O$, the points $A$ and $B$ of the unit circle are the images of $\dfrac{7\pi}6$ and $\dfrac{11\pi}6$. The unoriented angle $\widehat{AOB}$ will be given in radians, between 0 and $\pi$.
\begin{enumerate}
\item Give their exact coordinates. \hfill(1 pt)
\item Calculate $\vect{OA}\cdot\vect{OB}$. \hfill(1 pt)
\item Determine the unoriented angle AOB then AB. \hfill(2 pts)
\end{enumerate}
